% ============================================================================
% Fichier  : partie1/P01_C02_architecture_soc.tex
% Titre    : L'Architecture d'un SOC
% Auteur   : MINKA MI NGUIDJOÏ Thierry Emmanuel
% ============================================================================
% Sources : NIST SP 800-61r2 ; modèles de SOC (niveaux N1-N3) ; CIS Controls v8.
% ============================================================================

\chapter{L'Architecture d'un SOC}
\label{chap:soc}
\chaptermark{L'architecture d'un SOC}

\epigraph{%
    « Un SOC n'est pas une salle pleine d'écrans~;\\
    c'est un processus qui transforme des traces en décisions. »%
}{Principe du cours}

\niveauChapitre{Fondamental}{Chapitre~\ref{chap:pourquoi}.}

\begin{boxobjectifs}
\begin{itemize}
    \item Définir le centre d'opérations de sécurité (SOC) et sa mission.
    \item Décrire les niveaux d'analyse (N1, N2, N3) et leurs rôles.
    \item Comprendre les modèles d'organisation (interne, infogéré, hybride).
    \item Situer les fonctions chasse, ingénierie de détection, réponse.
    \item Dimensionner un SOC réaliste pour une PME.
\end{itemize}
\end{boxobjectifs}

\bigskip

% ============================================================================
\section*{Question directrice : de quelle équipe LiZbiZ a-t-elle besoin pour voir venir la prochaine attaque~?}
% ============================================================================

LiZbiZ n'a ni SOC ni analyste dédié~; M.~KOUNOU fait tout. Ce chapitre montre
ce qu'est un SOC, ses rôles et ses modèles, et comment en bâtir un
proportionné à une PME, seul ou avec le cabinet NYONG SÉCURITÉ.

\medskip
\textbf{Plan du chapitre.}
\begin{enumerate}
    \item La mission du SOC (section~\ref{sec:c02-mission}).
    \item Les niveaux d'analyse (section~\ref{sec:c02-niveaux}).
    \item Les modèles d'organisation (section~\ref{sec:c02-modeles}).
    \item Les fonctions transverses (section~\ref{sec:c02-fonctions}).
    \item Dimensionner pour une PME (section~\ref{sec:c02-pme}).
\end{enumerate}

\separateur

% ============================================================================
\section{La Mission du SOC}
\label{sec:c02-mission}
% ============================================================================

\begin{boxdef}
\termecle{SOC} (\textit{Security Operations Center})~: fonction, humaine et
technique, chargée de détecter, analyser et traiter les événements de
sécurité en continu, à partir des traces collectées par le \sigle{SIEM} et
les sondes.
\end{boxdef}

% ============================================================================
\section{Les Niveaux d'Analyse}
\label{sec:c02-niveaux}
% ============================================================================

\begin{tableaucours}{L{2.2cm}L{4.6cm}L{5.4cm}}{%
    \tableauHeader{Niveau} &
    \tableauHeader{Rôle} &
    \tableauHeader{Ce qu'il fait}
}
N1 & Surveillance, tri & Reçoit les alertes, écarte les faux positifs, qualifie, escalade \\
N2 & Investigation & Analyse l'alerte qualifiée, reconstitue, décide du traitement \\
N3 & Expertise, chasse & Forensique, chasse proactive, ingénierie de détection \\
\end{tableaucours}
\vspace{-0.8em}
\begin{center}{\small\itshape Tableau~2.1. Les niveaux d'un SOC}\end{center}

\begin{boxalert}
\textbf{Erreur fréquente.} Vouloir trois niveaux dès le premier jour.

\smallskip
\textbf{Règle~:} une PME commence souvent par un N1 outillé et un N2 à temps
partiel, l'expertise N3 étant externalisée. L'organisation suit la maturité,
pas l'inverse.
\end{boxalert}

% ============================================================================
\section{Les Modèles d'Organisation}
\label{sec:c02-modeles}
% ============================================================================

\begin{tableaucours}{L{3.0cm}L{4.8cm}L{4.4cm}}{%
    \tableauHeader{Modèle} &
    \tableauHeader{Principe} &
    \tableauHeader{Pour qui}
}
Interne & Équipe propre, 24/7 ou heures ouvrées & Grandes organisations \\
Infogéré (MSSP) & Un prestataire opère le SOC & PME sans équipe \\
Hybride & N1/N2 infogéré, N3 et décision internes & Cas fréquent et équilibré \\
\end{tableaucours}
\vspace{-0.8em}
\begin{center}{\small\itshape Tableau~2.2. Modèles d'organisation d'un SOC}\end{center}

% ============================================================================
\section{Les Fonctions Transverses}
\label{sec:c02-fonctions}
% ============================================================================

Au-delà des niveaux, un SOC mûr distingue trois fonctions~: la
\termecle{chasse} (rechercher activement ce qui n'a pas déclenché d'alerte,
Chapitre~\ref{chap:renseignement}), l'\termecle{ingénierie de détection}
(écrire et maintenir les règles, Chapitre~\ref{chap:regles}) et la
\termecle{réponse} (traiter l'incident, Partie~V). Elles se nourrissent l'une
l'autre~: chaque incident traité améliore les règles.

% ============================================================================
\section{Dimensionner pour une PME}
\label{sec:c02-pme}
% ============================================================================

\begin{boxscenario}{Opération BÉNOUÉ. Le SOC qui n'existait pas}
\textbf{Constatation.} LiZbiZ n'avait aucune fonction de détection~: pas de
SIEM, pas d'analyste, des journaux dispersés et non lus.

\medskip
\textbf{Analyse} (confiance élevée). Un modèle hybride convient~: le cabinet
NYONG SÉCURITÉ opère la surveillance (N1/N2) sur un SIEM mutualisé, LiZbiZ
garde la décision et le lien avec l'ANTIC. Le coût est compatible avec une
PME, et le délai de présence tombe de onze jours à quelques heures.
\end{boxscenario}

\begin{boxmetier}{RSSI}
\textbf{Choisir un modèle soutenable.} Le RSSI ne vise pas le SOC le plus
complet, mais celui que l'organisation peut financer et faire vivre dans la
durée. Un SOC interne sous-doté détecte moins bien qu'un infogéré bien piloté.
\end{boxmetier}

\separateur

\begin{boxresume}
\textbf{Ce qu'il faut retenir de ce chapitre~:}
\begin{itemize}
    \item \textbf{Le SOC} transforme des traces en décisions, en continu.
    \item \textbf{N1, N2, N3}~: tri, investigation, expertise.
    \item \textbf{Interne, infogéré, hybride}~: le modèle suit la maturité et
          le budget.
    \item \textbf{Chasse, ingénierie de détection, réponse}~: les fonctions
          transverses.
    \item \textbf{Pour une PME}, l'hybride est souvent le bon départ.
\end{itemize}
\end{boxresume}

\begin{boxinfo}
\textbf{Points de contrôle de l'élève.} Je sais~: définir un SOC~; décrire
N1/N2/N3~; comparer les modèles~; situer chasse, ingénierie, réponse~;
proposer un SOC pour une PME.
\end{boxinfo}

\bigskip

\begin{boxexo}[title={\ding{45}\quad Exercice~2.1. Les niveaux \niveaubase}]
Associez à N1, N2 ou N3~: écarter un faux positif~; faire du forensique~;
décider d'isoler un serveur.
\end{boxexo}

\begin{boxexo}[title={\ding{45}\quad Exercice~2.2. SOC ou exploitation \niveaubase}]
Pourquoi un SOC n'est-il pas « une salle pleine d'écrans »~?
\end{boxexo}

\begin{boxexo}[title={\ding{45}\quad Exercice~2.3. Modèles \niveaubase}]
Citez un avantage et une limite de chaque modèle (interne, infogéré,
hybride).
\end{boxexo}

\begin{boxexo}[title={\ding{45}\quad Exercice~2.4. Maturité \niveaumoyen}]
Pourquoi déconseille-t-on à une PME de monter d'emblée trois niveaux
internes~?
\end{boxexo}

\begin{boxexo}[title={\ding{45}\quad Exercice~2.5. Hybride \niveaumoyen}]
Dans le modèle hybride de LiZbiZ, qui garde la décision~? Pourquoi~?
\end{boxexo}

\begin{boxexo}[title={\ding{45}\quad Exercice~2.6. Fonctions \niveaumoyen}]
Expliquez comment un incident traité améliore les règles de détection.
\end{boxexo}

\begin{boxexo}[title={\ding{45}\quad Exercice~2.7. Escalade \niveaumoyen}]
Décrivez le parcours d'une alerte de N1 à N2, avec les critères d'escalade.
\end{boxexo}

\begin{boxexo}[title={\ding{45}\quad Exercice~2.8. Chasse \niveaumoyen}]
En quoi la chasse diffère-t-elle du traitement d'alerte~?
\end{boxexo}

\begin{boxexo}[title={\ding{45}\quad Exercice~2.9. Proposition \niveauexpert}]
Rédigez une proposition de SOC hybride pour LiZbiZ~: périmètre, rôles,
horaires, responsabilités partagées avec NYONG SÉCURITÉ.
\end{boxexo}

\begin{boxexo}[title={\ding{45}\quad Exercice~2.10. Convention de service \niveauexpert}]
Quels engagements de service (délais, couverture, rapports) inscririez-vous
au contrat avec un SOC infogéré~?
\end{boxexo}

\bigskip

\begin{boxinfo}
\textbf{Pour aller plus loin.}
\begin{itemize}
    \item NIST SP~800-61r2~\cite{nist80061r2}~; CIS Controls v8~\cite{cis2021controls}.
\end{itemize}
\end{boxinfo}

\begin{boxinfo}
\textbf{Transition.} Un SOC se pilote~: le Chapitre~\ref{chap:gouvernance}
traite de sa gouvernance et de ses indicateurs.
\end{boxinfo}
